\documentclass[a4paper,12pt]{report}

% --- PACCHETTI PRINCIPALI ---
\usepackage[T1]{fontenc}
\usepackage[utf8]{inputenc}
\usepackage[italian]{babel}
\usepackage{lmodern}
\usepackage{setspace}
\usepackage{hyperref}
\usepackage{graphicx}

\usepackage{titlesec}
\usepackage{tikz}
\usetikzlibrary{shapes, arrows, positioning}

\tikzset{
    crowfoot o-/.style={circle, draw, fill=white, minimum size=4pt, inner sep=0pt},
    crowfoot |-/.style={rectangle, draw, minimum size=4pt, inner sep=0pt},
    crowfoot ||-/.style={rectangle split, rectangle split parts=2, draw,
                         minimum width=4pt, minimum height=6pt, inner sep=0pt},
}

% --- LISTINGS (codice) ---
\usepackage{listings}
\usepackage{xcolor}

% Colori personalizzati
\definecolor{codebg}{RGB}{250, 244, 230}
\definecolor{codenumbers}{gray}{0.45}

\lstset{
    backgroundcolor=\color{codebg},
    basicstyle=\ttfamily\small,
    numbers=left,
    numberstyle=\scriptsize\color{codenumbers},
    numbersep=8pt,
    frame=none,
    showstringspaces=false,
    breaklines=true,
    breakatwhitespace=true,
    columns=fullflexible,
    keepspaces=true,
    tabsize=2,
}

% --- ALTRI PACCHETTI UTILI ---
\usepackage{geometry}
\usepackage{indentfirst}
\usepackage{booktabs}


\newcommand{\upref}[1]{\hyperref[#1]{\textsuperscript{$\uparrow$}}}

% --- IMPOSTAZIONI DI STILE ---
\geometry{a4paper, margin=2.5cm}
\onehalfspacing
\setlength{\parskip}{6pt}

% --- INTESTAZIONE TESI ---
\title{\textbf{SpaceParts}\\Piattaforma Tecnica per Componenti Aerospaziali e Robotici}
\author{Denis G. Boitor}
\date{Anno Accademico 2025}

\begin{document}


\maketitle  % <-- il titolo DEVE precedere l'indice

\clearpage

% --- INDICE DELLA TESINA ---
\tableofcontents

\clearpage

\chapter{Introduzione e descrizione dei dati}


\section{Specifiche introduttive}
\label{sec:introduzione}

La piattaforma \textbf{SpaceParts} nasce con l’obiettivo di offrire un ambiente completo, strutturato e facilmente navigabile per la gestione, la consultazione e l’acquisto di componenti destinati ad applicazioni aerospaziali, robotiche e di intelligenza artificiale. Essa si rivolge a un pubblico eterogeneo che include studenti, ricercatori, aziende, ingegneri e appassionati del settore, proponendosi come punto di incontro tra ambiti accademici, industriali e divulgativi.

L’accesso alla piattaforma SpaceParts avviene a partire da una pagina homepage visitabile da tutte le tipologie di utenza. Tale pagina fornisce una panoramica sintetica delle categorie merceologiche trattate, presenta i prodotti più rilevanti e mette in evidenza eventuali campagne promozionali globali in corso (come ad esempio eventi speciali dedicati a missioni spaziali, periodi commerciali o ricorrenze tecnologiche).

Il catalogo dei prodotti costituisce il cuore funzionale del sistema. Ogni componente è descritto mediante una scheda tecnica che ne riporta categoria, certificazione (ESA-grade, industrial, educational), anno di produzione, prezzo base, compatibilità con determinati sistemi, missioni a cui è associato, eventuali bonus crediti e sconti attivi. L’utente, anche qualora non autenticato, può sfogliare liberamente le schede tecniche e filtrare la visualizzazione tramite una serie di parametri selezionabili (categoria, fascia di prezzo, anno, certificazione, missione).

La piattaforma prevede la presenza di un sistema di \textbf{recensioni tecniche} che consente agli utenti registrati di contribuire al miglioramento delle informazioni sui prodotti acquistati. Il sistema delle recensioni consente agli utenti di commentare i prodotti acquistati. 
Ogni recensione contiene un voto (1–5) e un commento tecnico, e rappresenta un feedback puntuale sull’esperienza d’uso del componente. Le recensioni possono ricevere risposte da altri utenti, permettendo brevi discussioni tecniche correlate al prodotto.
A differenza del forum, non è previsto alcun sistema di voto delle recensioni: 
ciò garantisce che il sistema mantenga una funzione strettamente informativa, 
non sociale.
SpaceParts vuole supportare la ricerca e sviluppo tencologica, perciò presenta all'interno un \textbf{Forum Tecnico}, nella quale gli utenti possono porre domande, avviare discussioni e contribuire con risposte pertinenti a temi quali elettronica avanzata, propulsione, controllo, robotica e software di bordo.

Un elemento fondamentale del sistema è costituito dalla \textbf{reputazione utente}, la quale rappresenta un indicatore dell’affidabilità e della competenza dimostrata dal cliente all’interno della piattaforma.
Il cliente può partecipare alle discussioni votando i contenuti tramite un sistema di valutazione binario “utile / inutile”, rappresentato da una freccia verso l’alto e una verso il basso. Questo meccanismo permette di far emergere i contributi più pertinenti e influisce sulla reputazione dell’autore. La reputazione dell’utente viene inoltre 
influenzata dal numero di prodotti acquistati e dall’attività di recensione e dai giudizi ricevuti. Una reputazione più elevata permette agli utenti di accedere a particolari vantaggi, come sconti personalizzati o promozioni mirate.

La piattaforma gestisce inoltre un sistema di \textbf{crediti tecnici} che consente agli utenti autenticati di accumulare bonus in seguito all’acquisto di determinati prodotti o alla partecipazione a eventi promozionali. I crediti possono essere utilizzati come forma di sconto per acquisti futuri oppure ricaricati mediante apposita richiesta, soggetta all’approvazione dell’amministratore.

A livello di gestione amministrativa, il sistema include un modulo dedicato al controllo del catalogo e dei contenuti generati dagli utenti. All'interno di tale modulo, le figure con privilegi superiori possono aggiungere o modificare prodotti, applicare sconti, gestire le recensioni segnalate, approvare richieste di ricarica crediti e moderare i contenuti del forum.

SpaceParts suddivide chiaramente le tipologie di utenza, ciascuna dotata di una serie di funzionalità specifiche e non accessibili ad altre categorie. Le tipologie di utenti previste dalla piattaforma sono le seguenti:

\begin{itemize}
    \item \textbf{Visitatore} (utente non autenticato);
    \item \textbf{Cliente}, in grado di acquistare prodotti e contribuire ai contenuti;
    \item \textbf{Gestore del catalogo}, responsabile della gestione tecnica dei prodotti e delle promozioni;
    \item \textbf{Amministratore}, responsabile della gestione degli utenti, delle revisioni e dei contenuti.
\end{itemize}

Ogni ruolo dispone di funzionalità specifiche, alcune delle quali condivise e altre completamente esclusive. Nei paragrafi seguenti tali funzionalità verranno analizzate nel dettaglio, illustrando in modo chiaro e ordinato quali servizi sono disponibili per ciascuna tipologia di utenza.

\medskip

La sezione successiva è dedicata a un’analisi approfondita delle funzionalità specifiche per ogni tipologia di utente.

\section{Funzionalità specifiche}
\label{sec:funzionalita}

In questa sezione verranno approfondite nel dettaglio tutte le funzionalità
associate alle tipologie di utenza previste dalla piattaforma SpaceParts.

La piattaforma è progettata in modo da garantire che ciascun utente,
in relazione alla tipologia con la quale si autentica, possa accedere esclusivamente
alle funzionalità pertinenti, evitando così sovrapposizioni o accessi non autorizzati
a contenuti riservati.  
Le funzionalità sono suddivise in quattro gruppi principali:
Visitatore, Cliente, Gestore e Amministratore.


\subsection{Funzionalità Visitatore}
\label{sec:visitatore}

Un utente non autenticato viene definito \textbf{Visitatore}. Tale utenza
ha la possibilità di esplorare liberamente le parti pubbliche della piattaforma,
senza tuttavia poter interagire con le funzionalità avanzate riservate agli utenti registrati.
Le funzionalità disponibili sono le seguenti:

\begin{itemize}
    \item \textbf{Visualizzare il catalogo dei prodotti}.  
    La homepage permette all’utente di accedere direttamente all’intera lista
    dei componenti disponibili. Il catalogo è strutturato per categorie
    (sensori, elettronica di bordo, attuatori, propulsione elettrica,
    pannelli solari, moduli avionici, ecc.) e ogni voce può essere consultata
    liberamente.  
    Selezionando una categoria, l’utente viene reindirizzato a una pagina che
    elenca i prodotti corrispondenti, ciascuno corredato da un’immagine,
    un titolo e alcune informazioni tecniche preliminari.

    \item \textbf{Filtrare i prodotti}.  
    Il visitatore può applicare criteri di filtraggio basati su una o più proprietà:
    certificazione (ESA-grade, industrial, educational), anno di produzione,
    intervallo di prezzo, categoria tecnica, o missione spaziale di riferimento
    (es. “Artemis II”, “Gaia”, “ExoMars”).  
    Tale funzionalità permette di ridurre rapidamente la lista dei prodotti
    presentati, facilitando la ricerca di componenti specifici anche a chi
    non sia un esperto.

    \item \textbf{Accedere alla scheda tecnica di un prodotto}.  
    Ogni prodotto è associato a una scheda informativa dettagliata
    contenente descrizione, parametri operativi, requisiti di compatibilità,
    certificazioni e, in caso di promozioni attive, dettagli relativi a sconti
    e bonus crediti.  
    Pur non essendo autenticato, l’utente può esplorare tutte le schede tecniche.

    \item \textbf{Accedere alle sezioni pubbliche: Chi siamo, Missioni, Contatti}.  
    Queste sezioni offrono contenuti informativi sull’azienda o ente gestore
    del portale, sulle missioni aerospaziali supportate, e sulle modalità
    di contatto.

    \item \textbf{Accedere alle FAQ globali}.  
    Il visitatore può consultare un insieme di domande e risposte di carattere
    generale riguardanti i prodotti, la piattaforma o gli acquisti.

    \item \textbf{Accedere al Forum Tecnico (sola lettura)}.  
    L’utente può visualizzare discussioni tecniche attive, ma non può né
    creare contenuti né rispondere ai messaggi pubblicati.

    \item \textbf{Accedere alle pagine di login}.  
    Tramite apposito pulsante presente nella barra superiore,
    il visitatore una volta inserite le credenziali di accesso, sara' cura della piattaforma identificarlo come Cliente, Gestore o
    Amministratore.

    \item \textbf{Registrarsi alla piattaforma}.  
    Qualora l’utente desideri usufruire di funzionalità avanzate
    (come scrivere recensioni o effettuare acquisti),
    può accedere alla pagina di registrazione, nella quale dovrà fornire
    i dati richiesti (nome, cognome, email, username e password).
\end{itemize}

\subsection{Funzionalità Cliente}
\label{sec:cliente}
Un utente registrato ed autenticato come \textbf{Cliente} ottiene accesso a
funzionalità aggiuntive rispetto al Visitatore.  

Le funzionalità disponibili sono le seguenti:

\begin{itemize}
    \item \textbf{Visualizzare e modificare i propri dati personali}.  
    Dalla pagina principale dedicata all’utenza Cliente, è possibile accedere
    a un pannello riepilogativo contenente tutte le informazioni anagrafiche.
    L’utente può aggiornare alcuni dati, come l’email o la password.

    \item \textbf{Visualizzare lo storico degli acquisti}.  
    Una sezione dedicata elenca cronologicamente tutti gli ordini effettuati.
    Ogni ordine include la lista dei prodotti acquistati, le quantità,
    l’eventuale applicazione di crediti, il totale pagato e la data di acquisto.

    \item \textbf{Gestire il saldo crediti}.  
    All’interno del profilo utente è presente un indicatore che mostra
    l’attuale ammontare di crediti tecnici disponibili.
    L’utente può inoltre inviare una richiesta di ricarica,
    soggetta all’approvazione dell’amministratore.

    \item \textbf{Aggiungere prodotti al carrello}.  
    Visitando la scheda tecnica di un prodotto,
    il cliente può inserirlo nel carrello indicando, quando previsto,
    la quantità desiderata.

    \item \textbf{Effettuare acquisti}.  
    L’utente può procedere al checkout visualizzando il riepilogo dell’ordine:
    prodotti inclusi, sconti applicati, eventuali crediti utilizzabili,
    totale da pagare e bonus crediti che verranno maturati con l’acquisto.

    \item \textbf{Pubblicare recensioni tecniche}.  
    A seguito dell’acquisto di un prodotto,
    l’utente può lasciare una recensione valutando il componente
    tramite un punteggio da 1 a 5 stelle e un commento tecnico
    riguardante prestazioni, affidabilità, compatibilità e considerazioni pratiche.

    \item \textbf{Rispondere a recensioni esistenti}.  
    Una recensione può ospitare una breve discussione tecnica.
    L’utente può contribuire rispondendo con precisazioni, analisi tecniche
    o esperienze d’uso.

    \item \textbf{Valutare contenuti del Forum}.
    Il cliente può partecipare alle discussioni votando i contenuti tramite un sistema 
    di valutazione binario “utile / inutile”, rappresentato da una freccia verso l’alto 
    e una verso il basso. Questo meccanismo permette di far emergere i contributi più 
    pertinenti e influisce sulla reputazione dell’autore.

    \item \textbf{Accedere al Forum Tecnico}.  
    Il cliente può non solo visualizzare ma anche pubblicare nuovi post
    e rispondere a discussioni già esistenti.

    \item \textbf{Consultare le FAQ}.  
    L’utente può visualizzare domande e risposte pubblicate dagli Amministratori.
\end{itemize}

\subsection{Funzionalità Gestore}
\label{sec:gestore}

Il ruolo di \textbf{Gestore} rappresenta una figura intermedia tra il Cliente e l’Amministratore, con responsabilità dedicate alla gestione
del catalogo e alla moderazione dei contenuti tecnici.

Le funzionalità disponibili sono:

\begin{itemize}
    \item \textbf{Aggiungere nuovi prodotti}.  
    Il Gestore può compilare un modulo contenente nome, categoria,
    certificazione, prezzo, anno di produzione, bonus crediti e missioni associate,
    inserendo così un nuovo componente all’interno del file XML dedicato ai prodotti.

    \item \textbf{Modificare prodotti esistenti}.  
    Per ogni prodotto registrato, il Gestore può aggiornare informazioni
    come prezzo, certificazione, descrizione, anno, immagine e bonus.

    \item \textbf{Eliminare prodotti}.  
    Il sistema impedisce la rimozione di prodotti che risultino associati
    a ordini già effettuati, in modo da preservare la coerenza storica.
    In tutti gli altri casi il Gestore può rimuovere liberamente il prodotto.

    \item \textbf{Applicare sconti e bonus}.  
    Tramite un’interfaccia dedicata, il Gestore può definire sconti temporanei
    indicando percentuale, condizioni e durata, oppure assegnare bonus crediti.

    \item \textbf{Moderare discussioni e recensioni}.  
    Il Gestore può intervenire in caso di contenuti tecnici poco accurati,
    discussioni inappropriate o risposte non conformi allo scopo della piattaforma.
    È autorizzato a modificare o eliminare contenuti di bassa gravità.

\end{itemize}

\subsection{Funzionalità Amministratore}
\label{sec:admin}

L’\textbf{Amministratore} è la figura con il livello più elevato di privilegi,
ha la responsabilità dell’intera piattaforma e può intervenire in qualsiasi
sezione del sistema.

Le funzionalità disponibili sono:

\begin{itemize}
    \item \textbf{Gestione completa degli utenti}.  
    L’amministratore può visualizzare tutte le utenze registrate,
    modificarne i dati anagrafici, variare il ruolo, sospendere temporaneamente
    o bannare definitivamente un account.  
    Il ban temporaneo impedisce l’accesso a sezioni sensibili come recensioni
    e forum, mentre il ban definitivo impedisce qualsiasi accesso alla piattaforma.

    \item \textbf{Gestione delle recensioni e dei contenuti segnalati}.  
    L’amministratore può intervenire in caso di recensioni o discussioni
    ritenute inappropriate o tecnicamente scorrette,
    modificando o eliminando il contenuto.

    \item \textbf{Gestione delle richieste di ricarica crediti}.  
    Ogni richiesta effettuata da un cliente passa attraverso un processo di approvazione:
    l’amministratore può accettare o rifiutare la ricarica,
    registrando la decisione nei file XML dei crediti.

    \item \textbf{Modifica e gestione del Forum Tecnico}.  
    L’amministratore ha facoltà di riorganizzare discussioni,
    bloccare thread, eliminare contenuti e intervenire in caso di violazioni
    del regolamento.

    \item \textbf{Gestione delle FAQ}.  
    Può aggiungere nuove domande e risposte,
    aggiornare le FAQ esistenti o rimuovere contenuti obsoleti.

    \item \textbf{Gestione delle campagne promozionali globali}.  
    L’amministratore può definire eventi come ``Launch Week'' o
    ``Speciale Artemis II'', scegliendo durata, categoria interessata,
    percentuali di sconto e bonus associati.
    
\end{itemize}


\section{Strutture dati}
\label{sec:strutture-dati}

In questa sezione verranno illustrate nel dettaglio tutte le strutture dati che
la piattaforma SpaceParts utilizza per la memorizzazione e la gestione dei
contenuti:

\begin{itemize}
    \item una \textbf{tabella SQL} dedicata alla memorizzazione delle informazioni
          anagrafiche e dei ruoli degli utenti;
    \item una serie di \textbf{file XML} che rappresentano i contenuti dinamici
          del sistema (prodotti, recensioni, ordini, forum, FAQ, campagne, crediti).
\end{itemize}


\subsubsection{\texttt{Vincoli Relazionali}}

Di seguito si riportano alcuni vincoli logici per facilitare la comprensione.

\begin{figure}[h!]
\centering

\begin{tikzpicture}[
entity/.style={rectangle, draw, rounded corners, text width=4cm, minimum height=1cm, align=center},
rel/.style={diamond, draw, aspect=2, minimum width=1.8cm, align=center},
arrow/.style={-latex, thick}
]

% Entities
\node[entity] (prodotti) {PRODOTTI};
\node[entity, right=4cm of prodotti, yshift=1.5cm] (sconti) {SCONTI};
\node[entity, right=4cm of prodotti, yshift=-1.5cm] (recensioni) {RECENSIONI};

% Relations
\node[rel, midway, right=2cm of prodotti, yshift=1.5cm] (rel1) {1--N};
\node[rel, midway, right=2cm of prodotti, yshift=-1.5cm] (rel2) {1--N};

% Arrows
\draw[arrow] (prodotti) -- (rel1) -- (sconti);
\draw[arrow] (prodotti) -- (rel2) -- (recensioni);

\end{tikzpicture}

\caption{Diagramma relazionale tra Prodotti, Sconti e Recensioni.}
\label{fig:er}
\end{figure}


Il diagramma ER (Entity-Relationship) di Crow’s Foot rappresenta la struttura logica delle principali 
entità coinvolte nella piattaforma \textbf{SpaceParts}. L’entità \textbf{Prodotti} costituisce 
il nodo centrale del modello e presenta due relazioni di cardinalità uno-a-molti (1–N): 
una con l’entità \textbf{Sconti}, che descrive le promozioni applicabili a ciascun prodotto, 
e una con l’entità \textbf{Recensioni}, che raccoglie i contributi pubblicati dagli utenti. 
La presenza del simbolo a “piede di corvo” indica che un singolo prodotto può essere 
associato a molteplici sconti e a molteplici recensioni, mentre ogni sconto o recensione 
deve necessariamente riferirsi a un unico prodotto. Il modello riflette quindi una struttura 
dati coerente con i requisiti della piattaforma, garantendo estendibilità e integrità 

\subsection{Anagrafica utenti}
\label{sec: anagrafica-utenti}

La piattaforma prevede una tabella SQL denominata \textbf{utenti} contenente i dati essenziali per la gestione delle credenziali, del ruolo e degli aspetti amministrativi.


\subsection*{Utenti}

\begin{table}[h!]
\centering
\resizebox{\textwidth}{!}{%
\begin{tabular}{l l l l}
\toprule
\textbf{Campo} & \textbf{Tipo} & \textbf{Vincoli} & \textbf{Descrizione} \\
\midrule

\texttt{id} & INTEGER & \texttt{PRIMARY KEY, AUTO\_INCREMENT} & Identificativo univoco dell'utente \\

\texttt{username} & VARCHAR(15) & \texttt{UNIQUE, NOT NULL} & Nome utente scelto \\

\texttt{nome} & VARCHAR(15) & \texttt{NOT NULL} & Nome dell'utente \\

\texttt{cognome} & VARCHAR(15) & \texttt{NOT NULL} & Cognome dell'utente \\

\texttt{indirizzo} & VARCHAR(15) & \texttt{NOT NULL} & Indirizzo dell'utente \\

\texttt{email} & VARCHAR(15) & \texttt{UNIQUE, NOT NULL} & Indirizzo email \\

\texttt{password} & VARCHAR(15) & \texttt{NOT NULL} & Password criptata \\

\texttt{ruolo} & ENUM & \texttt{NOT NULL} & Ruolo dell'utente (admin/gestore/cliente) \\

\texttt{reputazione} & FLOAT & \texttt{DEFAULT 0} & Punteggio reputazione \\

\texttt{stato} & INTEGER & \texttt{DEFAULT 1} & Stato account (attivo/disattivo/sospeso) \\

\texttt{portafoglio} & FLOAT & \texttt{DEFAULT 0.00} & Importo disponibile \\

\texttt{crediti} & FLOAT & \texttt{DEFAULT 0} & Crediti tecnici accumulati \\

\texttt{data\_iscrizione} & DATE & \texttt{DEFAULT CURRENT\_DATE} & Data di registrazione \\
\bottomrule
\end{tabular}%
}
\end{table}


\subsection{Struttura dei file XML}
\label{sec:xml}

Ogni insieme di dati operativi viene gestito tramite un file XML dedicato, seguito dalla relativa grammatica definita tramite XSD, che sar`a poi utilizzata per la gestione dei dati
all’interno del codice dell’applicazione. I dati relativi agli utenti, come Amministratore, Gestore e Cliente verranno  gestiti tramite PHPMyAdmin, un'applicazione web che consente di amministrare un database MySQL.

% ------------------------------------------------------
% ------------------------------------------------------

\subsubsection{XML fondamentali}
\subsection*{Utenti}

Il file \texttt{utenti.xml} contiene esclusivamente informazioni non sensibili
relative ai profili degli utenti.

Le credenziali (username, password) e i ruoli amministrativi sono invece gestiti
tramite database MySQL, garantendo sicurezza e corretta separazione dei dati.

\subsection*{utenti.xml}
{\footnotesize
\begin{lstlisting}
<?xml version="1.0" encoding="UTF-8"?>
<utenti xmlns:xsi="http://www.w3.org/2001/XMLSchema-instance"
        xsi:noNamespaceSchemaLocation="utenti.xsd">

  <utente id="1">
    <nome>Admin</nome>
    <cognome>Admin</cognome>
    <telefono>3204012230</telefono>
    <indirizzo>Via Latina</indirizzo>
    <reputazione>0</reputazione>
    <portafoglio>0</portafoglio>
    <crediti>0</crediti>
    <data_iscrizione>2025-01-01</data_iscrizione>
  </utente>

</utenti>
\end{lstlisting}}

\subsection*{utenti.xsd}
{\footnotesize
\begin{lstlisting}
<?xml version="1.0" encoding="UTF-8"?>
<xs:schema xmlns:xs="http://www.w3.org/2001/XMLSchema"
           elementFormDefault="qualified">

  <xs:element name="utenti">
    <xs:complexType>
      <xs:sequence>
        <xs:element name="utente" maxOccurs="unbounded">
          <xs:complexType>
            <xs:sequence>
              <xs:element name="nome" type="xs:string"/>
              <xs:element name="cognome" type="xs:string"/>
              <xs:element name="telefono" type="xs:string"/>
              <xs:element name="indirizzo" type="xs:string"/>
              <xs:element name="reputazione" type="xs:integer"/>
              <xs:element name="portafoglio" type="xs:decimal"/>
              <xs:element name="crediti" type="xs:decimal"/>
              <xs:element name="data_iscrizione" type="xs:date"/>
            </sequence>
            <xs:attribute name="id" type="xs:integer" use="required"/>
          </xs:complexType>
        </xs:element>
      </sequence>
    </xs:complexType>
  </xs:element>

</xs:schema}
\end{lstlisting}}

% ------------------------------------------------------

\subsection*{Prodotti}
\label{sec:catalogo}
Il file XML dedicato ai \textbf{prodotti} memorizza l’intero catalogo tecnico della piattaforma SpaceParts.  
Ogni elemento \texttt{<prodotto>} contiene un identificativo univoco e i principali attributi utilizzati dal sistema:
nome, categoria, descrizione, prezzo, bonus crediti, certificazione, anno di produzione, missione associata e percorso dell’immagine rappresentativa.  


\subsubsection*{Prodotti.XML}

\begin{lstlisting}
<?xml version="1.0" encoding="UTF-8"?>
<prodotti>
  <prodotto id="1">
    <nome>Unità di Controllo Avionica A320-X</nome>
    <categoria>Avionica</categoria>
    <descrizione>
      Modulo avionico ad alta affidabilità con processore
      rad-hard certificato ESA.
    </descrizione>
    <prezzo>2950.00</prezzo>
    <bonus_crediti>30</bonus_crediti>
    <certificazione>ESA-grade</certificazione>
    <anno_produzione>2024</anno_produzione>
    <missione>Artemis II</missione>
    <immagine>a320x_control_unit.jpg</immagine>
  </prodotto>
  <prodotto id="2">
    <nome>Sensore Termico T900 UltraSense</nome>
    <categoria>Sensori</categoria>
    <descrizione>
      Sensore termico con precisione a 0.01°C,
      ideale per sistemi robotici e controllo satelliti.
    </descrizione>
    <prezzo>690.00</prezzo>
    <bonus_crediti>7</bonus_crediti>
    <certificazione>Industrial</certificazione>
    <anno_produzione>2023</anno_produzione>
    <missione>Earth Observation</missione>
    <immagine>t900_thermal.jpg</immagine>
  </prodotto>
        .
        .
        .
</prodotti>
\end{lstlisting}

\subsubsection*{Prodotti.XSD}

{\footnotesize
\begin{lstlisting}
<?xml version="1.0" encoding="UTF-8"?>
<xs:schema xmlns:xs="http://www.w3.org/2001/XMLSchema" elementFormDefault="qualified">
  <xs:element name="prodotti">
    <xs:complexType>
      <xs:sequence>
        <xs:element name="prodotto" maxOccurs="unbounded">
          <xs:complexType>
            <xs:sequence>
              <xs:element name="nome" type="xs:string"/>
              <xs:element name="categoria" type="xs:string"/>
              <xs:element name="descrizione" type="xs:string"/>
              <xs:element name="prezzo" type="xs:decimal"/>
              <xs:element name="bonus_crediti" type="xs:integer"/>
              <xs:element name="certificazione" type="certificazioneType"/>
              <xs:element name="anno_produzione" type="xs:gYear"/>
              <xs:element name="missione" type="xs:string"/>
              <xs:element name="immagine" type="xs:string"/>
            </xs:sequence>
            <xs:attribute name="id" type="xs:integer" use="required"/>
          </xs:complexType>
        </xs:element>
      </xs:sequence>
    </xs:complexType>
  </xs:element>
  <xs:simpleType name="certificazioneType">
    <xs:restriction base="xs:string">
      <xs:enumeration value="ESA-grade"/>
      <xs:enumeration value="Industrial"/>
      <xs:enumeration value="Space-certified"/>
      <xs:enumeration value="Educational"/>
    </xs:restriction>
  </xs:simpleType>
</xs:schema>
\end{lstlisting}
}

% ------------------------------------------------------

\subsection*{Sconti}
\label{sec:xml-sconti}

Il file \texttt{sconti.xml} contiene l'elenco degli sconti applicabili ai prodotti
presenti nel catalogo SpaceParts. Ogni elemento \texttt{<sconto>} definisce una
promozione specifica associata a un singolo prodotto e include:

\begin{itemize}
    \item \textbf{id\_prodotto}: il prodotto a cui lo sconto si applica;
    \item \textbf{percentuale}: riduzione percentuale del prezzo;
    \item \textbf{condizione}: eventuale requisito per attivare lo sconto 
          (evento promozionale, reputazione minima, ecc.);
    \item \textbf{data\_inizio} e \textbf{data\_fine}: intervallo temporale di validità.
\end{itemize}

Questo file supporta il sistema delle promozioni personalizzate e globali,
integrandosi con le campagne e con il catalogo prodotti.

\subsection*{sconti.xml}
{\footnotesize
\begin{lstlisting}
<?xml version="1.0" encoding="UTF-8"?>
<sconti xmlns:xsi="http://www.w3.org/2001/XMLSchema-instance"
        xsi:noNamespaceSchemaLocation="sconti.xsd">

  <sconto id="1">
    <id_prodotto>1</id_prodotto>
    <percentuale>15</percentuale>
    <condizione>Evento Artemis Week</condizione>
    <data_inizio>2025-02-01</data_inizio>
    <data_fine>2025-02-10</data_fine>
  </sconto>

</sconti>
\end{lstlisting}}

\subsection*{sconti.xsd}
{\footnotesize
\begin{lstlisting}
<?xml version="1.0" encoding="UTF-8"?>
<xs:schema xmlns:xs="http://www.w3.org/2001/XMLSchema"
           elementFormDefault="qualified">

  <xs:element name="sconti">
    <xs:complexType>
      <xs:sequence>
        <xs:element name="sconto" maxOccurs="unbounded">
          <xs:complexType>
            <xs:sequence>
              <xs:element name="id_prodotto" type="xs:integer"/>
              <xs:element name="percentuale" type="xs:decimal"/>
              <xs:element name="condizione" type="xs:string"/>
              <xs:element name="data_inizio" type="xs:date"/>
              <xs:element name="data_fine" type="xs:date"/>
            </xs:sequence>
            <xs:attribute name="id" type="xs:integer" use="required"/>
          </xs:complexType>
        </xs:element>
      </xs:sequence>
    </xs:complexType>
  </xs:element>

</xs:schema}
\end{lstlisting}}


% ------------------------------------------------------
% ------------------------------------------------------

\subsection{Sistema recensioni}
\subsection*{Recensioni}
\label{sec:xml-recensioni}

Il file \texttt{recensioni.xml} contiene tutte le recensioni tecniche pubblicate
dagli utenti a seguito dell’acquisto di un prodotto.  
Ogni elemento \texttt{<recensione>} include:

\begin{itemize}
    \item \textbf{id\_prodotto}: riferimento al prodotto recensito;
    \item \textbf{id\_utente}: autore della recensione;
    \item \textbf{valutazione}: punteggio assegnato (1–5);
    \item \textbf{commento}: analisi tecnica o impressioni di utilizzo;
    \item \textbf{data}: data di pubblicazione.
\end{itemize}

Le risposte alle recensioni non sono incluse in questo file,
ma sono archiviate separatamente in \texttt{risposte\_recens.xml},
così da mantenere una struttura più modulare e scalabile.

\subsection*{recensioni.xml}
{\footnotesize
\begin{lstlisting}
<?xml version="1.0" encoding="UTF-8"?>
<recensioni xmlns:xsi="http://www.w3.org/2001/XMLSchema-instance"
            xsi:noNamespaceSchemaLocation="recensioni.xsd">

  <recensione id="1">
    <id_prodotto>1</id_prodotto>
    <id_utente>2</id_utente>
    <valutazione>5</valutazione>
    <commento>Componente estremamente affidabile.</commento>
    <data>2025-01-05</data>
  </recensione>

</recensioni>
\end{lstlisting}}

\subsection*{recensioni.xsd}
{\footnotesize
\begin{lstlisting}
<?xml version="1.0" encoding="UTF-8"?>
<xs:schema xmlns:xs="http://www.w3.org/2001/XMLSchema"
           elementFormDefault="qualified">

  <xs:element name="recensioni">
    <xs:complexType>
      <xs:sequence>
        <xs:element name="recensione" maxOccurs="unbounded">
          <xs:complexType>
            <xs:sequence>
              <xs:element name="id_prodotto" type="xs:integer"/>
              <xs:element name="id_utente" type="xs:integer"/>
              <xs:element name="valutazione" type="xs:integer"/>
              <xs:element name="commento" type="xs:string"/>
              <xs:element name="data" type="xs:date"/>
            </xs:sequence>
            <xs:attribute name="id" type="xs:integer" use="required"/>
          </xs:complexType>
        </xs:element>
      </xs:sequence>
    </xs:complexType>
  </xs:element>

</xs:schema}
\end{lstlisting}}

% ------------------------------------------------------

\subsection*{Risposte alle Recensioni}
\label{sec:xml-risposte}

Il file \texttt{risposte\_recens.xml} contiene le risposte associate alle recensioni tecniche.
Ogni risposta si riferisce a una recensione esistente.

\subsection*{risposte\_recens.xml}

\begin{lstlisting}
<?xml version="1.0" encoding="UTF-8"?>
<risposte_recens xmlns:xsi="http://www.w3.org/2001/XMLSchema-instance"
                 xsi:noNamespaceSchemaLocation="risposte_recens.xsd">

  <risposta id="1">
    <id_recensione>1</id_recensione>
    <autore>4</autore>
    <testo>Confermo le ottime prestazioni del modulo avionico.</testo>
    <data>2025-01-07</data>
  </risposta>

</risposte_recens>
\end{lstlisting}

\subsection*{risposte\_recens.xsd}

\begin{lstlisting}
<?xml version="1.0" encoding="UTF-8"?>
<xs:schema xmlns:xs="http://www.w3.org/2001/XMLSchema"
           elementFormDefault="qualified">

  <xs:element name="risposte_recens">
    <xs:complexType>
      <xs:sequence>
        <xs:element name="risposta" maxOccurs="unbounded">
          <xs:complexType>
            <xs:sequence>
              <xs:element name="id_recensione" type="xs:integer"/>
              <xs:element name="autore" type="xs:integer"/>
              <xs:element name="testo" type="xs:string"/>
              <xs:element name="data" type="xs:date"/>
            </xs:sequence>
            <xs:attribute name="id" type="xs:integer" use="required"/>
          </xs:complexType>
        </xs:element>
      </xs:sequence>
    </xs:complexType>
  </xs:element>

</xs:schema}
\end{lstlisting}

% ------------------------------------------------------
% ------------------------------------------------------

\subsection{Sistema interazione comunità}
\subsection*{Forum}
\label{sec:xml-forum}

Il file \texttt{forum.xml} archivia le discussioni tecniche pubblicate dagli utenti.
Ogni elemento \texttt{<discussione>} contiene:

\begin{itemize}
    \item \textbf{titolo}: argomento della discussione;
    \item \textbf{autore}: utente che ha avviato il thread;
    \item \textbf{data} e \textbf{testo}: contenuto principale;
    \item \textbf{voti_up/down}: tieni traccia dei voti utili e/o inutili
    \item risposte interne, rappresentate da elementi \texttt{<risposta>}.
\end{itemize}

Questo file riflette la struttura lineare del forum, utile per la consultazione
tecnica collettiva.

\subsection*{forum.xml}
{\footnotesize
\begin{lstlisting}
<?xml version="1.0" encoding="UTF-8"?>
<forum xmlns:xsi="http://www.w3.org/2001/XMLSchema-instance"
       xsi:noNamespaceSchemaLocation="forum.xsd">

  <discussione id="1">
    <titolo>Problemi di propulsione ionica</titolo>
    <autore>2</autore>
    <data>2025-01-03</data>
    <testo>Qualcuno ha avuto problemi con il modello Helios-I2?</testo>
    <voti_up>12</voti_up>
    <voti_down>1</voti_down>
      <risposta id="1">
     <autore>3</autore>
     <data>2025-01-04</data>
     <testo>Funziona correttamente in condizioni nominali.</testo>
     <voti_up>12</voti_up>
     <voti_down>1</voti_down>
      </risposta>

  </discussione>

</forum>
\end{lstlisting}}

\subsection*{forum.xsd}
{\footnotesize
\begin{lstlisting}
<?xml version="1.0" encoding="UTF-8"?>
<xs:schema xmlns:xs="http://www.w3.org/2001/XMLSchema"
           elementFormDefault="qualified">
  <xs:element name="forum">
    <xs:complexType>
      <xs:sequence>
        <xs:element name="discussione" maxOccurs="unbounded">
          <xs:complexType>
            <xs:sequence>
              <xs:element name="titolo" type="xs:string"/>
              <xs:element name="autore" type="xs:integer"/>
              <xs:element name="data" type="xs:date"/>
              <xs:element name="testo" type="xs:string"/>
              <xs:element name="voti_up" type="xs:integer"/>
              <xs:element name="voti_down" type="xs:integer"/>

              <xs:element name="risposta" maxOccurs="unbounded" minOccurs="0">
                <xs:complexType>
                  <xs:sequence>
                    <xs:element name="autore" type="xs:integer"/>
                    <xs:element name="data" type="xs:date"/>
                    <xs:element name="testo" type="xs:string"/>
                    <xs:element name="voti_up" type="xs:integer"/>
                    <xs:element name="voti_down" type="xs:integer"/>

                  </xs:sequence>
                  <xs:attribute name="id" type="xs:integer" use="required"/>
                </xs:complexType>
              </xs:element>
            </xs:sequence>
            <xs:attribute name="id" type="xs:integer" use="required"/>
          </xs:complexType>
        </xs:element>
      </xs:sequence>
    </xs:complexType>
  </xs:element>
</xs:schema}
\end{lstlisting}
}

% ------------------------------------------------------

\subsection*{FAQ}
\label{sec:xml-faq}

Il file \texttt{faq.xml} raccoglie le domande più frequenti pubblicate dagli amministratori
per assistere gli utenti nella navigazione della piattaforma.
Ogni voce comprende:

\begin{itemize}
    \item \textbf{quesito}: la domanda formulata;
    \item \textbf{risposta}: la spiegazione associata.
\end{itemize}

Tale struttura agevola la gestione centralizzata delle informazioni.

\subsection*{faq.xml}
{\footnotesize
\begin{lstlisting}
<?xml version="1.0" encoding="UTF-8"?>
<faq xmlns:xsi="http://www.w3.org/2001/XMLSchema-instance"
     xsi:noNamespaceSchemaLocation="faq.xsd">

  <domanda id="1">
    <quesito>Come posso acquistare un componente?</quesito>
    <risposta>Accedi alla piattaforma e utilizza il carrello.</risposta>
  </domanda>

</faq>
\end{lstlisting}}

\subsection*{faq.xsd}
{\footnotesize
\begin{lstlisting}
<?xml version="1.0" encoding="UTF-8"?>
<xs:schema xmlns:xs="http://www.w3.org/2001/XMLSchema"
           elementFormDefault="qualified">

  <xs:element name="faq">
    <xs:complexType>
      <xs:sequence>
        <xs:element name="domanda" maxOccurs="unbounded">
          <xs:complexType>
            <xs:sequence>
              <xs:element name="quesito" type="xs:string"/>
              <xs:element name="risposta" type="xs:string"/>
            </xs:sequence>
            <xs:attribute name="id" type="xs:integer" use="required"/>
          </xs:complexType>
        </xs:element>
      </xs:sequence>
    </xs:complexType>
  </xs:element>

</xs:schema}
\end{lstlisting}
}

% ------------------------------------------------------
% ------------------------------------------------------

\subsection{Sistema economico}
\label{sec:xml-economia}
\subsection*{crediti}


Il file \texttt{crediti.xml} registra tutte le operazioni relative al saldo crediti
degli utenti: bonus ottenuti tramite acquisti, ricariche approvate dagli amministratori
o rettifiche di sistema.  
Ogni movimento è identificato da:

\begin{itemize}
    \item \textbf{id\_utente}
    \item \textbf{importo}
    \item \textbf{tipo} (bonus, ricarica, utilizzo)
    \item \textbf{data}
\end{itemize}

Questo file garantisce la tracciabilità delle operazioni economiche.


\subsection*{crediti.xml}

{\footnotesize
\begin{lstlisting}
<?xml version="1.0" encoding="UTF-8"?>
<crediti xmlns:xmi="http://www.w3.org/2001/XMLSchema-instance"
         xsi:noNamespaceSchemaLocation="crediti.xsd">

  <movimento id="1">
    <id_utente>2</id_utente>
    <importo>50</importo>
    <data>2025-01-02</data>
    <tipo>ricarica</tipo>
  </movimento>

</crediti>
\end{lstlisting}
}

\subsection*{crediti.xsd}

{\footnotesize
\begin{lstlisting}
<?xml version="1.0" encoding="UTF-8"?>
<xs:schema xmlns:xs="http://www.w3.org/2001/XMLSchema"
           elementFormDefault="qualified">

  <xs:element name="crediti">
    <xs:complexType>
      <xs:sequence>

        <xs:element name="movimento" maxOccurs="unbounded">
          <xs:complexType>
            <xs:sequence>
              <xs:element name="id_utente" type="xs:integer"/>
              <xs:element name="importo" type="xs:decimal"/>
              <xs:element name="data" type="xs:date"/>
              <xs:element name="tipo" type="xs:string"/>
            </xs:sequence>
            <xs:attribute name="id" type="xs:integer" use="required"/>
          </xs:complexType>
        </xs:element>

      </xs:sequence>
    </xs:complexType>
  </xs:element>

</xs:schema}
\end{lstlisting}
}

% ------------------------------------------------------

\subsection{Richieste di Ricarica Crediti}

Il file \texttt{rich\_crediti.xml} contiene tutte le richieste di ricarica inviate dagli utenti.
Ogni richiesta deve essere approvata o rifiutata da un amministratore.

\subsection*{rich\_crediti.xml}

\begin{lstlisting}
<?xml version="1.0" encoding="UTF-8"?>
<rich_crediti xmlns:xsi="http://www.w3.org/2001/XMLSchema-instance"
              xsi:noNamespaceSchemaLocation="rich_crediti.xsd">

  <richiesta id="1">
    <id_utente>3</id_utente>
    <importo>100</importo>
    <stato>in_attesa</stato>
    <data>2025-01-06</data>
  </richiesta>

</rich_crediti>
\end{lstlisting}

\subsection*{rich\_crediti.xsd}

\begin{lstlisting}
<?xml version="1.0" encoding="UTF-8"?>
<xs:schema xmlns:xs="http://www.w3.org/2001/XMLSchema"
           elementFormDefault="qualified">

  <xs:element name="rich_crediti">
    <xs:complexType>
      <xs:sequence>

        <xs:element name="richiesta" maxOccurs="unbounded">
          <xs:complexType>
            <xs:sequence>
              <xs:element name="id_utente" type="xs:integer"/>
              <xs:element name="importo" type="xs:decimal"/>
              <!-- in_attesa, approvata, rifiutata -->
              <xs:element name="stato" type="xs:string"/>
              <xs:element name="data" type="xs:date"/>
            </xs:sequence>
            <xs:attribute name="id" type="xs:integer" use="required"/>
          </xs:complexType>
        </xs:element>

      </xs:sequence>
    </xs:complexType>
  </xs:element>

</xs:schema}
\end{lstlisting}

% ------------------------------------------------------

\subsection*{Campagne promozionali}

Le campagne promozionali consentono all’amministratore di applicare sconti temporanei a insiemi mirati di prodotti.
Per rendere il sistema più flessibile, ogni campagna può essere associata a una o più categorie del catalogo (ad es. Avionica, Sensori, Propulsione).
L’elemento <categorie> permette infatti di specificare l’ambito di applicazione della promozione senza intervenire sui singoli prodotti, migliorando modularità e scalabilità del sistema.
Se nessuna categoria viene indicata, la campagna è considerata globale e si applica all’intero catalogo.

\subsection*{Confronto tra Campagne promozionali e Sconti}

\begin{itemize}
    \item \textbf{Campagne promozionali}: operano a livello di \textit{categoria},
    applicando una riduzione di prezzo a gruppi di prodotti omogenei
    (ad esempio tutti i componenti di Avionica). Si tratta di iniziative globali,
    pianificate dall’amministratore e legate a periodi o eventi specifici.

    \item \textbf{Sconti}: agiscono a livello di \textit{singolo prodotto}.
    Ogni sconto è associato a un preciso identificativo di prodotto e può prevedere
    condizioni particolari (evento promozionale, reputazione minima, promozioni mirate, ecc.).
\end{itemize}

\noindent In sintesi, le campagne hanno un effetto \textit{esteso} e tematico,
mentre gli sconti hanno un effetto \textit{puntuale} e mirato.


\subsection*{campagne.xml}

{\footnotesize
\begin{lstlisting}
<?xml version="1.0" encoding="UTF-8"?>
<campagne xmlns:xsi="http://www.w3.org/2001/XMLSchema-instance"
          xsi:noNamespaceSchemaLocation="campagne.xsd">

  <campagna id="1">
    <nome>Avionics Boost</nome>
    <descrizione>Sconti limitati sui moduli avionici certificati.</descrizione>
    <data_inizio>2025-03-01</data_inizio>
    <data_fine>2025-03-10</data_fine>
    <categorie>
      <categoria>Avionica</categoria>
      <categoria>Propulsione</categoria>
    </categorie>
  </campagna>

  <campagna id="2">
    <nome>Spring Sensors Week</nome>
    <descrizione>-20% su tutti i sensori ambientali.</descrizione>
    <data_inizio>2025-04-01</data_inizio>
    <data_fine>2025-04-07</data_fine>
    <categorie>
      <categoria>Sensori</categoria>
    </categorie>
  </campagna>
        .
        .
        .
</campagne>
\end{lstlisting}
}

\subsection*{campagne.xsd}

{\footnotesize
\begin{lstlisting}
<?xml version="1.0" encoding="UTF-8"?>
<xs:schema xmlns:xs="http://www.w3.org/2001/XMLSchema"
           elementFormDefault="qualified">

  <xs:element name="campagne">
    <xs:complexType>
      <xs:sequence>

        <xs:element name="campagna" maxOccurs="unbounded">
          <xs:complexType>
            <xs:sequence>

              <xs:element name="nome" type="xs:string"/>
              <xs:element name="descrizione" type="xs:string"/>
              <xs:element name="data_inizio" type="xs:date"/>
              <xs:element name="data_fine" type="xs:date"/>

              <!-- elenco di categorie collegate alla campagna -->
              <xs:element name="categorie" minOccurs="0">
                <xs:complexType>
                  <xs:sequence>
                    <xs:element name="categoria" type="xs:string"                       maxOccurs="unbounded"/>
                  </xs:sequence>
                </xs:complexType>
              </xs:element>

            </xs:sequence>

            <xs:attribute name="id" type="xs:integer" use="required"/>
          </xs:complexType>
        </xs:element>

      </xs:sequence>
    </xs:complexType>
  </xs:element>

</xs:schema>

\end{lstlisting}
}


% ------------------------------------------------------
% ------------------------------------------------------


\subsection{Sistema di acquisto}
\subsection*{Carrello}

Il file \texttt{carrelli.xml} contiene la rappresentazione persistente del carrello
associato agli utenti. Ogni elemento \texttt{<carello>} identifica il carrello di
un singolo utente, mentre gli elementi \texttt{<item>} elencano i prodotti
attualmente presenti nel carrello con la relativa quantità.

\subsection*{carelli.xml}
\begin{lstlisting}
<?xml version="1.0" encoding="UTF-8"?>
<carrelli xmlns:xsi="http://www.w3.org/2001/XMLSchema-instance"
          xsi:noNamespaceSchemaLocation="carrelli.xsd">

  <carrello id_utente="1">
    <item id_prodotto="1" quantita="1"/>
    <item id_prodotto="4" quantita="2"/>
  </carrello>

</carrelli>

\end{lstlisting}

\subsection*{carello.xsd}
\begin{lstlisting}
<?xml version="1.0" encoding="UTF-8"?>
<xs:schema xmlns:xs="http://www.w3.org/2001/XMLSchema"
           elementFormDefault="qualified">

  <xs:element name="carrelli">
    <xs:complexType>
      <xs:sequence>

        <xs:element name="carrello" maxOccurs="unbounded">
          <xs:complexType>
            <xs:sequence>

              <xs:element name="item" maxOccurs="unbounded">
                <xs:complexType>
                  <xs:attribute name="id_prodotto" type="xs:integer" use="required"/>
                  <xs:attribute name="quantita" type="xs:integer" use="required"/>
                </xs:complexType>
              </xs:element>

            </xs:sequence>
            <xs:attribute name="id_utente" type="xs:integer" use="required"/>
          </xs:complexType>
        </xs:element>

      </xs:sequence>
    </xs:complexType>
  </xs:element>

</xs:schema>

\end{lstlisting}

% ------------------------------------------------------

\subsection*{Ordini}

Il file \texttt{ordini.xml} descrive ogni transazione effettuata dagli utenti,
includendo prodotti acquistati, quantità, data e totale dell'ordine.

\subsection*{ordini.xml}

{\footnotesize
\begin{lstlisting}
<?xml version="1.0" encoding="UTF-8"?>
<ordini xmlns:xsi="http://www.w3.org/2001/XMLSchema-instance"
        xsi:noNamespaceSchemaLocation="ordini.xsd">

  <ordine id="1">
    <id_utente>3</id_utente>
    <data>2025-01-04</data>
    <totale>3800.00</totale>

    <prodotto id="1" quantita="1"/>
    <prodotto id="2" quantita="2"/>
  </ordine>

</ordini>
\end{lstlisting}
}

\subsection*{ordini.xsd}

{\footnotesize
\begin{lstlisting}
<?xml version="1.0" encoding="UTF-8"?>
<xs:schema xmlns:xs="http://www.w3.org/2001/XMLSchema"
           elementFormDefault="qualified">

  <xs:element name="ordini">
    <xs:complexType>
      <xs:sequence>

        <xs:element name="ordine" maxOccurs="unbounded">
          <xs:complexType>
            <xs:sequence>

              <xs:element name="id_utente" type="xs:integer"/>
              <xs:element name="data" type="xs:date"/>
              <xs:element name="totale" type="xs:decimal"/>

              <xs:element name="prodotto" maxOccurs="unbounded">
                <xs:complexType>
                  <xs:attribute name="id" type="xs:integer" use="required"/>
                  <xs:attribute name="quantita" type="xs:integer" use="required"/>
                </xs:complexType>
              </xs:element>

            </xs:sequence>
            <xs:attribute name="id" type="xs:integer" use="required"/>
          </xs:complexType>
        </xs:element>

      </xs:sequence>
    </xs:complexType>
  </xs:element>

</xs:schema}
\end{lstlisting}
}

% ------------------------------------------------------

\subsection*{Storico Acquisti}

Il file \texttt{storico\_acq.xml} conserva la cronologia completa degli acquisti
effettuati dagli utenti, indipendentemente dallo stato del carrello o delle promozioni.

\subsection*{storico\_acq.xml}

\begin{lstlisting}
<?xml version="1.0" encoding="UTF-8"?>
<storico_acq xmlns:xsi="http://www.w3.org/2001/XMLSchema-instance"
             xsi:noNamespaceSchemaLocation="storico_acq.xsd">

  <acquisto id="1">
    <id_utente>3</id_utente>
    <data>2025-01-04</data>
    <totale>3800.00</totale>
    <prodotto id="1" quantita="1"/>
    <prodotto id="2" quantita="2"/>
  </acquisto>

</storico_acq>
\end{lstlisting}

\subsection*{storico\_acq.xsd}

\begin{lstlisting}
<?xml version="1.0" encoding="UTF-8"?>
<xs:schema xmlns:xs="http://www.w3.org/2001/XMLSchema"
           elementFormDefault="qualified">

  <xs:element name="storico_acq">
    <xs:complexType>
      <xs:sequence>

        <xs:element name="acquisto" maxOccurs="unbounded">
          <xs:complexType>
            <xs:sequence>
              <xs:element name="id_utente" type="xs:integer"/>
              <xs:element name="data" type="xs:date"/>
              <xs:element name="totale" type="xs:decimal"/>
              <xs:element name="prodotto" maxOccurs="unbounded">
                <xs:complexType>
                  <xs:attribute name="id" type="xs:integer" use="required"/>
                  <xs:attribute name="quantita" type="xs:integer" use="required"/>
                </xs:complexType>
              </xs:element>
            </xs:sequence>
            <xs:attribute name="id" type="xs:integer" use="required"/>
          </xs:complexType>
        </xs:element>

      </xs:sequence>
    </xs:complexType>
  </xs:element>

</xs:schema}
\end{lstlisting}

% ------------------------------------------------------
% ------------------------------------------------------


\subsection{Controllo qualità}
\label{sec:xml-controllo}


\subsection*{Segnalazioni}

Il file \texttt{segnalazioni.xml} elenca le segnalazioni inviate dagli utenti
per errori tecnici, problemi nei prodotti o comportamenti inappropriati nel forum.

\subsection*{segnalazioni.xml}

\begin{lstlisting}
<?xml version="1.0" encoding="UTF-8"?>
<segnalazioni xmlns:xsi="http://www.w3.org/2001/XMLSchema-instance"
              xsi:noNamespaceSchemaLocation="segnalazioni.xsd">

  <segnalazione id="1">
    <id_utente>5</id_utente>
    <categoria>forum</categoria>
    <testo>Messaggio non conforme alle linee guida.</testo>
    <data>2025-01-08</data>
  </segnalazione>

</segnalazioni>
\end{lstlisting}

\subsection*{segnalazioni.xsd}

\begin{lstlisting}
<?xml version="1.0" encoding="UTF-8"?>
<xs:schema xmlns:xs="http://www.w3.org/2001/XMLSchema"
           elementFormDefault="qualified">

  <xs:element name="segnalazioni">
    <xs:complexType>
      <xs:sequence>

        <xs:element name="segnalazione" maxOccurs="unbounded">
          <xs:complexType>
            <xs:sequence>
              <xs:element name="id_utente" type="xs:integer"/>
              <xs:element name="categoria" type="xs:string"/>
              <xs:element name="testo" type="xs:string"/>
              <xs:element name="data" type="xs:date"/>
            </xs:sequence>
            <xs:attribute name="id" type="xs:integer" use="required"/>
          </xs:complexType>
        </xs:element>

      </xs:sequence>
    </xs:complexType>
  </xs:element>

</xs:schema}
\end{lstlisting}


% ------------------------------------------------------
% ------------------------------------------------------


\section{Analisi di alcuni elementi e attributi dei linguaggi XML}
\label{sec:analisi-xml}

In questa sezione vengono analizzati alcuni elementi e attributi scelti a 
campione dai linguaggi XML definiti nel capitolo precedente.  
Per ciascun elemento viene giustificata l’esistenza all’interno del modello 
informativo, facendo riferimento alle specifiche funzionali esposte nelle sezioni
precedenti.

% --------------------------------------------------------
\subsection{File \texttt{prodotti.xml}}

\subsubsection*{Attributo \texttt{id}}
L’attributo \texttt{id} identifica in maniera univoca ogni prodotto del catalogo.  
La sua presenza è necessaria poiché il prodotto è la risorsa principale 
richiamata da numerosi altri documenti XML.  
Come descritto nella Sezione~\ref{sec:catalogo}, la piattaforma deve 
mantenere riferimenti consistenti ai prodotti durante visualizzazione, ricerca, 
acquisto e gestione amministrativa.  
L’identificativo garantisce quindi integrità referenziale tra i vari moduli del sistema.

\subsubsection*{Elemento \texttt{categoria}}
L’elemento \texttt{categoria}, presente nella sezione ~\ref{sec:catalogo}, permette di classificare i prodotti in gruppi tecnici (ad esempio avionica, sensori, robotica), dove la ricerca deve supportare la selezione per categorie.  
L’assenza di tale elemento renderebbe impossibile strutturare il catalogo secondo criteri funzionali. Inoltre non permetterebbe l'appliczione di eventuali sconti o promozioni globali


% --------------------------------------------------------
\subsection{File \texttt{sconti.xml}}

\subsubsection*{Elemento \texttt{id\_prodotto}}
L’elemento collega ogni sconto ad uno specifico prodotto del catalogo.  
Questa relazione è coerente con il modello relazionale introdotto nella 
Figura~\ref{fig:er}.  
Questo elemento, presente nella Sezione~\ref{sec:xml-sconti}, dà la possibilità di applicare promozioni selettive basate su prodotti o categorie.  
L’elemento \texttt{id\_prodotto} è quindi indispensabile per collegare il prezzo 
dinamico al prodotto corretto.

\subsubsection*{Elemento \texttt{percentuale}}
Rappresenta la riduzione di prezzo applicata al prodotto. Una volta definito lo sconto da applicare all'i-esimo utente (in base alla reputazione e/o crediti, campagne promozionali)
Visitabile nella Sezione~\ref{sec:xml-sconti}, questo valore è quindi necessario per aggiornare dinamicamente il prezzo visualizzato.

% --------------------------------------------------------
% --------------------------------------------------------
\subsection{File \texttt{recensioni.xml}}

\subsubsection*{Elemento \texttt{valutazione}}
L’elemento \texttt{valutazione}, un intero compreso tra 1 e 5, permette agli utenti di esprimere 
un giudizio sintetico sulle prestazioni e sull’affidabilità del prodotto recensito.  
Come descritto nella Sezione~\ref{sec:xml-recensioni}, la piattaforma richiede 
che ogni recensione includa un punteggio numerico standardizzato, così da favorire 
il confronto immediato tra componenti appartenenti a categorie diverse.  
In assenza di tale elemento, il sistema di feedback risulterebbe meno strutturato, 
poiché le opinioni rimarrebbero esclusivamente testuali e meno confrontabili.

\subsubsection*{Elementi \texttt{voti\_up} e \texttt{voti\_down}}
Questi elementi registrano rispettivamente il numero di valutazioni positive e negative 
ricevute da una discussione o da una risposta nel forum.  
Come illustrato nella Sezione introduttiva~\ref{sec:introduzione}, il sistema 
di voto permette di far emergere i contributi più rilevanti, migliorando la qualità 
complessiva delle interazioni tecniche. Essi contribuiscono inoltre alla reputazione 
dell’autore dei contenuti.


% --------------------------------------------------------
\subsection{File \texttt{carrello.xml}}

\subsubsection*{Attributo \texttt{id\_utente}}
L’attributo \texttt{id\_utente} (vedi Sez.~\ref{sec:xml-economia}) identifica in modo univoco il proprietario del carrello e 
permette di mantenere lo stato del carrello tra sessioni diverse.  
il sistema deve garantire che ogni utente disponga di un carrello persistente, capace di conservare prodotti e quantità 
anche dopo la disconnessione.  
Senza tale attributo, la piattaforma non potrebbe associare il contenuto del carrello allo 
specifico utente autenticato.

\subsubsection*{Attributi \texttt{id\_prodotto} e \texttt{quantita}}
Gli attributi \texttt{id\_prodotto} e \texttt{quantita}, presenti negli elementi \texttt{item},
descrivono la composizione effettiva del carrello: quali prodotti sono stati aggiunti e in quale quantità.  
La procedura di finalizzazione 
dell’ordine si basa proprio sulla riconciliazione tra prodotti nel carrello, sconti attivi e 
quantità selezionate.  
La mancanza di tali attributi impedirebbe di ricostruire il riepilogo dell’acquisto,
rendendo impossibile generare lo scontrino finale o aggiornare correttamente il saldo crediti.

% --------------------------------------------------------
\subsection{File \texttt{segnalazioni.xml}}

\subsubsection*{Elemento \texttt{categoria}}
L’elemento \texttt{categoria} (vedi Sez.~\ref{sec:xml-controllo}) distingue il tipo di contenuto oggetto della segnalazione
(ad esempio: prodotto, recensione, risposta, discussione del forum).  
Il sistema di moderazione richiede una classificazione precisa per permettere agli amministratori di intervenire 
in modo mirato a seconda della natura del contenuto segnalato.  
Questo elemento è quindi fondamentale per indirizzare rapidamente la segnalazione al modulo corretto del sistema (catalogo, recensioni, forum, ecc.).

\subsubsection*{Elemento \texttt{testo}}
L’elemento \texttt{testo} contiene la descrizione fornita dall’utente riguardo al problema 
riscontrato.  
Gli utenti devono essere messi in condizione di motivare chiaramente la propria segnalazione, 
in modo da assistere i moderatori nell’analisi del caso.  
Questo campo consente di raccogliere dettagli contestuali indispensabili 
per distinguere un semplice malinteso da una violazione effettiva delle norme della piattaforma.


\end{document}
